% 场景一 S1-C 改进算法轨迹（子图 b）
\begin{tikzpicture}[scale=0.85, transform shape]
  \fill[bglight] (0,-1.5) rectangle (12,1.5);
  \draw[deepblue, very thick] (0,-1.5) -- (12,-1.5);
  \draw[deepblue, very thick] (0, 1.5) -- (12, 1.5);
  \draw[deepblue, dashed, thin] (0,0) -- (12,0);

  \fill[obsstatic] (2.2,-0.2) rectangle (2.7,0.35);
  \fill[obsstatic] (4.0,-0.35) rectangle (4.5,0.2);
  \fill[obsstatic] (5.8,-0.1) rectangle (6.3,0.45);
  \fill[obsstatic] (7.6,-0.4) rectangle (8.1,0.15);
  \fill[obsstatic] (9.4,-0.2) rectangle (9.9,0.35);

  \fill[lightgreen!30] (0,-1.0) rectangle (12,-0.7);
  \draw[deeppink, very thick]
    plot[smooth, tension=0.5] coordinates {(0.3,0) (1.5,-0.6) (3.0,-0.85) (4.5,-0.85) (6.0,-0.85) (7.5,-0.85) (9.0,-0.85) (10.5,-0.6) (11.7,-0.2)};

  \fill[egopt] (0.3,-0.2) rectangle (0.9,0.2);
  \fill[deeppink] (11.7,-0.2) circle (3pt);

  \node[lightgreen!40!black, font=\scriptsize, anchor=west] at (0.1,-0.85) {最优通行带};
\end{tikzpicture}
